\section{附录：符号与术语表}\label{sec:appendix}

\begin{table}[htbp]
    \centering
    \caption{本文使用的主要符号}
    \label{tab:notation}
    \begin{tabular}{ll}
        \toprule
        符号 & 含义 \\
        \midrule
        $\pi_2(x)$ & $\#\{p \le x : p,\ p+2 \text{ 均为素数}\}$ \\
        $C_2$ & 孪生素数常数 $\prod_{p>2}(1 - (p-1)^{-2}) \approx 0.66016$ \\
        $B_2$ & 布朗常数（孪生素数倒数和） \\
        $p_n$ & 第 $n$ 个素数 \\
        $\Lambda(n)$ & von Mangoldt 函数 \\
        $\psi(x)$ & 切比雪夫函数 $\sum_{n \le x} \Lambda(n)$ \\
        $\theta$ & 素数在算术级数中的分布水平 \\
        $H$ & 有界素数间隔的界 \\
        $S(\mathcal{A}, z)$ & 集合 $\mathcal{A}$ 中不被小于 $z$ 的素数整除的元素个数 \\
        $\lambda_d$ & 塞尔伯格筛权 \\
        $P_r$ & 至多 $r$ 个素因子（计重数）的整数 \\
        EH & Elliott--Halberstam 猜想 \\
        \bottomrule
    \end{tabular}
\end{table}

\section*{术语对照}
\begin{itemize}
    \item 孪生素数（twin primes）：相差为 $2$ 的素数对；
    \item 筛法（sieve method）：通过容斥估计不被一组素数整除的整数个数的方法；
    \item 奇偶性壁垒（parity problem）：经典筛法无法区分素因子个数奇偶性的根本限制；
    \item 分布水平（level of distribution）：素数误差在算术级数模数上的平均可控范围；
    \item 可允许 $k$ 元组（admissible $k$-tuple）：不被任何单个素数整除筛尽的整数组。
\end{itemize}
